% =====================================================================
%  二、问题分析
% =====================================================================
\section{问题分析}

\subsection{总体思路}
四个子问题层层递进：问题一建立最基本的示向度交会几何定位模型；问题二将其
扩展为检测点的最优实验设计；问题三在前述基础上，面对多个未知源与严格的时间
预算，形成信息驱动的大范围搜索与清除策略；问题四进一步处理定向干扰源带来的
方向不确定性与诊断困难。\textbf{【待填充：补充总体技术路线，可用一张流程图
（图~1）概括。】}

\subsection{问题一分析}
问题一的本质是\textbf{计算几何问题}：每条示向度带有 $\pm \delta$ 的误差，形成
一条扇形射线束，两条射线束的交叠区域即定位区域；区域直径可由凸包直径
（旋转卡壳算法）求得，而“直径圆能否覆盖该区域”则归结为最小覆盖圆问题。
\textbf{【待填充：细化分析，说明“直径圆可能无法覆盖细长区域”的直觉。】}

\subsection{问题二分析}
问题二具有\textbf{最优实验设计}的思想：示向度误差 $\delta$ 在距离 $d$ 处产生的
横向误差约为 $d\tan\delta$，故“离目标近”与“两条示向度交会角接近 $90^\circ$”
之间存在权衡，可用 Fisher 信息矩阵 / GDOP / 误差椭圆加以形式化。
\textbf{【待填充：给出第二检测点选择的优化准则与候选区域的刻画方法。】}

\subsection{问题三分析}
问题三是本题的主战场，受\textbf{严格现实时间预算}（$20$ 分钟）约束。仅靠
“边扫边定”的低效策略难以完成，需要利用“频道数量有限且互不相同”这一结构，
先粗扫建立“频道—射线束”台账，再对交会区域收敛的频道精确定位，本质上是一个
\textbf{信息驱动的旅行商问题（TSP with active sensing）}。
\textbf{【待填充：展开时间预算核算与整体策略框架。】}

\subsection{问题四分析}
问题四叠加了定向干扰源（方向未知），使“测不到信号”出现了多种截然不同的
原因（无该频道、超出有效接收半径、处于定向覆盖角之外），诊断逻辑复杂度显著
上升。\textbf{【待填充：分析定向源对定位与清除流程的影响与应对思路。】}
